\documentclass[a4paper,10pt,twoside]{article}

\usepackage[top=1in, bottom=1in, left=1in, right=1in]{geometry}
\usepackage[utf8]{inputenc}
\usepackage[spanish,es-ucroman,es-noquoting]{babel}
\usepackage{setspace}
\usepackage{fancyhdr}
\usepackage{lastpage}
\usepackage{amsmath}
\usepackage{amsfonts}
\usepackage{verbatim}
\usepackage{graphicx}
\usepackage{float}
\usepackage{algpseudocode}
\usepackage{enumitem} % Provee macro \setlist
\usepackage[toc, page]{appendix}


%%%%%%%%%% Configuración de Fancyhdr - Inicio %%%%%%%%%%
\pagestyle{fancy}
\thispagestyle{fancy}
\lhead{Trabajo Práctico 2 · Sistemas Operativos}
\rhead{Delgado · Lovisolo · Petaccio}
\renewcommand{\footrulewidth}{0.4pt}
\cfoot{\thepage /\pageref{LastPage}}

\fancypagestyle{caratula} {
   \fancyhf{}
   \cfoot{\thepage /\pageref{LastPage}}
   \renewcommand{\headrulewidth}{0pt}
   \renewcommand{\footrulewidth}{0pt}
}
%%%%%%%%%% Configuración de Fancyhdr - Fin %%%%%%%%%%


%%%%%%%%%% Configuración de Algorithmic - Inicio %%%%%%%%%%
% Entorno propio para customizar la presentación del pseudocódigo
\newenvironment{pseudo}[1][]{%
    \vspace{0.5em}%
    \begin{algorithmic}%
}
{%
    \end{algorithmic}%
    \vspace{0.5em}%
}

% Llamada a función para usar así: \Fn{Foo}{bar, baz}.
% Produce \textsc{Foo}$(bar, baz)$.
\newcommand{\Fn}[2]{\textsc{#1}$(#2)$}

% Cláusula return para usar así: \PReturn foo
\newcommand{\PReturn}[1]{\textbf{return} $#1$}

% Cláusula break
\newcommand{\Break}{\textbf{break}}

% Operadores lógicos
\newcommand{\PAnd}{\textbf{and} }
\newcommand{\POr}{\textbf{or} }
\newcommand{\PNot}{\textbf{not} }

% Booleanos
\newcommand{\PTrue}{\textnormal{TRUE} }
\newcommand{\PFalse}{\textnormal{FALSE} }

% Null
\newcommand{\Null}{\textnormal{NULL} }

% Conectivo 'in' para usar así: \ForAll{$foo$ \In $bar$}
\newcommand{\In}{\textbf{in} }

% Conectivo 'to' para usar así: \For{$i = 1$ \In $n$}
\newcommand{\To}{\textbf{to} }

% Complejidades
\newcommand{\Ode}[1]{\hfill $O(#1)$}
%%%%%%%%%% Configuración de Algorithmic - Fin %%%%%%%%%%


%%%%%%%%%% Configuración de Appendix - Inicio %%%%%%%%%%
% Asigna la traducción de la palabra 'Appendices'.
\renewcommand{\appendixtocname}{Apéndices}
\renewcommand{\appendixpagename}{Apéndices}
%%%%%%%%%% Configuración de Appendix - Fin %%%%%%%%%%


%%%%%%%%%% Miscelánea - Inicio %%%%%%%%%%
% Evita que el documento se estire verticalmente para ocupar el espacio vacío
% en cada página.
\raggedbottom

% Deshabilita sangría en la primer línea de un párrafo.
\setlength{\parindent}{0em}

% Separación entre párrafos.
\setlength{\parskip}{0.5em}

% Separación entre elementos de listas.
\setlist{itemsep=0.5em}
%%%%%%%%%% Miscelánea - Fin %%%%%%%%%%


\begin{document}


%%%%%%%%%%%%%%%%%%%%%%%%%%%%%%%%%%%%%%%%%%%%%%%%%%%%%%%%%%%%%%%%%%%%%%%%%%%%%%%
%% Carátula                                                                  %%
%%%%%%%%%%%%%%%%%%%%%%%%%%%%%%%%%%%%%%%%%%%%%%%%%%%%%%%%%%%%%%%%%%%%%%%%%%%%%%%


\thispagestyle{caratula}

\begin{center}

\includegraphics[height=2cm]{DC.png} 
\hfill
\includegraphics[height=2cm]{UBA.jpg} 

\vspace{2cm}

Departamento de Computación,\\
Facultad de Ciencias Exactas y Naturales,\\
Universidad de Buenos Aires

\vspace{4cm}

\begin{Huge}
Trabajo Práctico 3
\end{Huge}

\vspace{0.5cm}

\begin{Large}
Sistemas Operativos
\end{Large}

\vspace{1cm}

Segundo Cuatrimestre de 2013

\vspace{4cm}

\begin{tabular}{|c|c|c|}
\hline
Apellido y Nombre & LU & E-mail\\
\hline
Alejandro Nahuel Delgado & 601/11 & nahueldelgado@gmail.com\\
Leandro Lovisolo         & 645/11 & leandro@leandro.me\\
Lautaro José Petaccio    & 443/11 & lausuper@gmail.com\\
\hline
\end{tabular}

\end{center}

\newpage


%%%%%%%%%%%%%%%%%%%%%%%%%%%%%%%%%%%%%%%%%%%%%%%%%%%%%%%%%%%%%%%%%%%%%%%%%%%%%%%
%% Índice                                                                    %%
%%%%%%%%%%%%%%%%%%%%%%%%%%%%%%%%%%%%%%%%%%%%%%%%%%%%%%%%%%%%%%%%%%%%%%%%%%%%%%%


\tableofcontents

\newpage


%%%%%%%%%%%%%%%%%%%%%%%%%%%%%%%%%%%%%%%%%%%%%%%%%%%%%%%%%%%%%%%%%%%%%%%%%%%%%%%
%% Implementación con MongoDB                                                %%
%%%%%%%%%%%%%%%%%%%%%%%%%%%%%%%%%%%%%%%%%%%%%%%%%%%%%%%%%%%%%%%%%%%%%%%%%%%%%%%

\section{Implementación con MongoDB}

\newpage

%%%%%%%%%%%%%%%%%%%%%%%%%%%%%%%%%%%%%%%%%%%%%%%%%%%%%%%%%%%%%%%%%%%%%%%%%%%%%%%
%% Análisis de escalabilidad                                                 %%
%%%%%%%%%%%%%%%%%%%%%%%%%%%%%%%%%%%%%%%%%%%%%%%%%%%%%%%%%%%%%%%%%%%%%%%%%%%%%%%

\section{Análisis de escalabilidad}
Antes de comenzar con el análisis en sí de las soluciones disponibles en el mercado, veamos algunas definiciones necesarias para comprender y clasificar las opciones.

\subsection{Introducción}
Cloud Computing es un término muy discutido durante los últimos años, y la cantidad de servicios que lo proveen aumenta año tras año. Muchas empresas han incursionado en este modelo de negocio, y muchas más han sido las que han aprovechado su potencial. Aún así, no hay ningún tipo de estandarización de los servicios ni de la definición de los mismos conceptos. En principio vamos a tratar de establecer algunas guías y definiciones, para luego poder centrarnos en nuestras necesidades particulares alrededor de mongoDB.

\subsection{Definición de \emph{Cloud Computing}}
Una nube consiste en un conjunto de computadoras conectadas y virtualizadas que son dinámicamente proveídas y presentadas como uno o más recursos computacionales unificados. Los clientes se conectan a esta nube para poder hacer uso del servicio que se provee.

En general, los servicios de cloud computing comprenden lo siguiente:
\begin{itemize}
	\item Abstracción de la infraestructura y software subyacente para ser ofrecidos como un servicio.
	\item Mantenido sobre una infraestructura escalable y flexible.
	\item Ofrece el servicio a demanda y garantías de calidad de servicio.
	\item Se abona por el uso de recursos computacionales sin necesidad de que los usuarios se comprometan por adelantado.
	\item Compartido y multiusuario.
	\item Accesible a través de internet desde cualquier dispositivo.
\end{itemize}
Cloud computing podría definirse como un modelo para permitir acceso conveniente, a demada y a través de la red a un conjunto compartido y configurable de recursos computacionales que pueden ser rápidamente proveídos y lanzados con un mínimo esfuerzo de administración o interacción con el proveedor.

Existen tres grandes modelos de servicio en los que generalmente se diferencian a los servicios de cloud computing: Software-as-a-Service (SaaS), Platform-as-a-Serive (PaaS e Infrastructure-as-a-Service (IaaS).

\subsection{Alternativas de \emph{IaaS} y \emph{PaaS}}

\subsection{De la nube al hardware propio}


\end{document}